\section{Qualitätssicherungsphase}
\label{sec:Qualitätssicherungsphase}
Um, nach Abschluss der Realisierungsphase, die Vollständigkeit der Umsetzung gegenüber den umzusetzenden Anforderungen und die
Funktionstüchtigkeit der Erweiterung sicherzustellen, werden zahlreiche Tests durchgeführt.


\subsection{Definition of Done}
\label{sub:Definition of Done}
Das Definition of Done (DoD) ist eine Art Vereinbarung, die definiert wann eine Aufgabe, wie die Erweiterung, als fertig anzusehen
ist. In diesem Fall wurde innerhalb des Entwicklerteams eine solche Vereinbarung getroffen, dabei steht die Sicherstellung der
Funktionalität und die Wahrung der Softwarequalität im Vordergrund. Beim Sichten des Programmcodes wird nicht nur das eigene
Verständnis dessen geschärft, sondern auch verbessert und schlanker gemacht. Frei nach dem Prinzip des Pfadfinders, wie es im
Rotem Grad der Clean Code Developer Initiative\footnote{Vgl.
\url{http://www.clean-code-developer.de/Roter-Grad.ashx\#Praktiken_10} (Aufruf: 02.04.2015).} beschrieben ist:

"`Hinterlasse einen Ort immer in einem besseren Zustand als du ihn vorgefunden hast."'

Da zeitgleich komplexere Abschnitte des Programmes kommentiert werden, wird dieser auch verständlicher. Durch die sinnvollen
Kommentare, inklusive JavaDoc, kann auf eine externe Entwicklerdokumentation verzichtet werden.

\begingroup
\setlength\intextsep{0pt}
\begin{wraptable}[4]{r}[-24pt]{0\textwidth}

\begin{tabular}{rl|c}

\hline
\rowcolor{head} \multicolumn{1}{c}{\textbf{Nr.}}   & \multicolumn{1}{c|}{\textbf{Testfall}} &
\multicolumn{1}{|c}{\textbf{Auswertung}} \\
\hline
1     & Consistent\linebreak \Fachbegriff{Repository} & \checkmark\\
\hline
\rowcolor{odd}2     & Task statistics validated & \checkmark\\
\hline
3     & Connectivity & \checkmark\\
\hline
\rowcolor{odd}4     & Incoming\linebreak event & \checkmark\\
\hline
5     & Data transfer & \checkmark\\
\hline
\rowcolor{odd}6     & Frontend event handling & \checkmark\\
\hline
7     & User scoping & \checkmark\\
\hline
\rowcolor{odd}8     & Accessibility & \checkmark\\
\hline
9     & Browser\linebreak compatibility & \checkmark\\

\end{tabular}

\caption{Testergebnisse}
\label{tab:Testergebnisse}

\end{wraptable}



% \begin{tabular}{rp{2.5cm}|p{5.5cm}|c}
% 
% \hline
% \multicolumn{1}{c}{\textbf{Nr.}}   & \multicolumn{1}{c|}{\textbf{Testfall}} & \multicolumn{1}{c}{\textbf{Erwartetes Ergebnis}} &
% \multicolumn{1}{|c}{\textbf{Auswertung}} \\
% \hline
% 1     & Consistent\linebreak \Fachbegriff{Repository} & Daten sind vorhanden und können nur vom \Fachbegriff{Repository}
% modifiziert werden & \checkmark\\
% \hline
% \rowcolor{odd}2     & Task statistics validated & Korrektes Ergebnis der Tasks-Statistik wird geliefert & \checkmark\\
% \hline
% 3     & Connectivity & Client-Server Verbindung besteht (\acs{XML}HttpRequest auf "pending") & \checkmark\\
% \hline
% \rowcolor{odd}4     & Event\linebreak incoming & Im Frontend wird die JavaScript-Methode \Methode{onMessage} angesprochenn &
% \checkmark\\
% \hline
% 5     & Data transfer & Event-Daten sind im \Abkuerzung{JSON}-Format & \checkmark\\
% \hline
% \rowcolor{odd}6     & Frontend event handling & Task-Event sorgt für Aktualisierung der \Abkuerzung{UI} & \checkmark\\
% \hline
% 7     & User scoping & Nur gewünschte Benutzer werden benachrichtigt & \checkmark\\
% \hline
% \rowcolor{odd}8     & Accessibility  & Status-Anzeigen sind barrierefrei Nutzbar & \checkmark\\
% \hline
% 9     & Browser\linebreak compatibility & Das Design funktioniert in allen Browser angemessen  & \checkmark\\
% 
% \end{tabular}

\subsection{Testergebnisse}
\label{sub:Testergebnisse}
Das Resultat der ausgeführten Tests wird in der \Tabelle{tab:Testergebnisse} festgehalten.  In \TabelleV{tab:Testfälle} sind die
durchgeführten Tests beschrieben.
